% ============================================================================
% Fichier  : partie4/P04_C09_acces_habilitations.tex
% Titre    : Accès et Habilitations
% Auteur   : MINKA MI NGUIDJOÏ Thierry Emmanuel
% Version  : 0.1 -- novembre 2026
% ============================================================================
% Sources : INF3226 C18, C21 (accès) ; COSO principe 11 ; ISO/IEC 27002:2022
% (gestion des accès, résumé) ; ACSI LAB-02 (matrice des droits).
% Faits MINÉQUIR établis ici (LAB-03) : 47 comptes actifs au 23/03/2026 ;
% 9 comptes d'agents partis ; 3 comptes génériques (saisie01, daf, admin) ;
% 5 comptes LIQ + ORD (dont mbarga.j) ; 2 comptes FRS + ORD (mbarga.j,
% nkolo.r) ; profil ADM : atangana.p et admin (générique) ; mot de passe
% sans expiration ; dernière revue des droits : aucune.
% ============================================================================

\chapter{Accès et Habilitations}
\label{chap:acces}

\epigraph{%
    « Un droit d'accès est une signature que l'on peut prêter. »%
}{Principe du cours}

\niveauChapitre{Intermédiaire}{Chapitres~\ref{chap:eval-sci}
et~\ref{chap:preuve}.}

\begin{boxobjectifs}
\begin{itemize}
    \item Décrire le cycle de vie d'un compte et les contrôles attendus à
          chaque étape.
    \item Repérer les comptes à risque~: partis, génériques, à privilèges,
          inactifs.
    \item Croiser la liste des comptes avec le fichier du personnel.
    \item Analyser une matrice de profils et identifier les conflits de
          séparation des tâches.
    \item Évaluer la revue périodique des habilitations.
    \item Coter les contrôles d'accès et classer les déficiences.
\end{itemize}
\end{boxobjectifs}

\bigskip

% ============================================================================
\section*{Pourquoi ce chapitre}
% ============================================================================

Les accès sont le premier des contrôles généraux informatiques, et le
plus important~: s'ils ne sont pas maîtrisés, aucun autre contrôle n'est
fiable (Chapitre~\ref{chap:sci}). Le double paiement du prologue est
d'abord une défaillance des accès. Ce chapitre donne la grille complète
pour les évaluer.

\separateur

% ============================================================================
\section{Le cycle de vie d'un compte}
\label{sec:c09-cycle}
% ============================================================================

\begin{tableaucours}{L{2.4cm}L{4.8cm}L{5.0cm}}{%
    \tableauHeader{Étape} &
    \tableauHeader{Contrôle attendu} &
    \tableauHeader{Preuve à demander}
}
Création & Demande écrite, validée par le chef de service, pour un agent
identifié & Formulaires de demande, liste des créations \\
Attribution des profils & Profils conformes à la fiche de poste~; pas
de combinaison interdite & Matrice des profils, demandes \\
Modification & Nouvelle demande validée lors d'un changement de poste &
Historique des profils \\
Revue périodique & Le propriétaire de l'application valide la liste des
comptes et des profils & Revues signées, avec suites \\
Désactivation & À la date d'effet du départ ou de la mutation & Journal
des désactivations, rapproché du fichier du personnel \\
\end{tableaucours}
\vspace{-0.8em}
\begin{center}{\small\itshape Tableau~9.1. Le cycle de vie d'un
compte}\end{center}

\begin{boxdef}
\termecle{Propriétaire d'application}\index{propriétaire d'application}~:
responsable métier qui décide des droits dans une application (pour
SIDÉPENSE, le directeur des affaires financières). La cellule
informatique \emph{exécute} les décisions du propriétaire~; elle ne les
prend pas.
\end{boxdef}

\separateur

% ============================================================================
\section{Les comptes à risque}
\label{sec:c09-risque}
% ============================================================================

\begin{tableaucours}{L{2.8cm}L{4.4cm}L{5.0cm}}{%
    \tableauHeader{Type} &
    \tableauHeader{Risque} &
    \tableauHeader{Test}
}
Compte d'un agent parti & Utilisé par l'ancien titulaire ou par un
tiers & Croiser avec le fichier du personnel \\
Compte générique (\texttt{saisie01}) & Aucune action n'est imputable à
une personne & Lister les comptes sans titulaire nominatif \\
Compte à privilèges (ADM) & Peut tout faire, y compris effacer ses traces
& Lister les détenteurs, vérifier la journalisation de leurs actions \\
Compte inactif & Oublié, donc non surveillé & Dernière connexion de plus
de 90 jours \\
\end{tableaucours}
\vspace{-0.8em}
\begin{center}{\small\itshape Tableau~9.2. Les comptes à risque}\end{center}

\begin{boxoutils}{Croiser les comptes avec le fichier du personnel}
Deux feuilles dans le même classeur~: \texttt{COMPTES} (colonne A
matricule) et \texttt{PERSONNEL} (colonne A matricule, colonne D
situation). Dans \texttt{COMPTES}~:
\begin{lstlisting}
E2 : =SIERREUR(RECHERCHEV(A2;PERSONNEL!$A:$D;4;FAUX);"ABSENT")
\end{lstlisting}
Un filtre sur \texttt{E} donne les comptes d'agents mutés ou partis
(situation autre que «~en poste~») et les comptes sans titulaire au
fichier (\texttt{ABSENT}), qui sont souvent des comptes génériques.
\end{boxoutils}

\begin{boxscenario}{Dossier MINÉQUIR. Les 47 comptes}
\textbf{Constatation.} Au 23~mars~2026, SIDÉPENSE compte 47 comptes
actifs. Le croisement avec le fichier du personnel donne 9 comptes
d'agents mutés ou partis et 3 comptes sans titulaire~:
\texttt{saisie01}, \texttt{daf} et \texttt{admin}. Le profil ADM est
détenu par \texttt{atangana.p} et par le compte générique
\texttt{admin}. Les mots de passe n'expirent jamais.

\medskip
\textbf{Analyse} (confiance élevée). Le compte \texttt{admin} est le plus
grave~: il peut tout faire, et rien de ce qu'il fait n'est imputable à une
personne. Le journal du prologue ne serait d'aucune aide si les
opérations du 28~décembre avaient été faites sous ce compte.
\end{boxscenario}

\separateur

% ============================================================================
\section{La matrice des profils}
\label{sec:c09-matrice}
% ============================================================================

L'analyse part de la matrice extraite par l'auditeur (un compte par ligne,
un profil par colonne) et des paires interdites du Tableau~2.6.

\begin{boxoutils}{Repérer les conflits dans la matrice}
Colonnes \texttt{B} à \texttt{F}~: ENG, LIQ, ORD, FRS, ADM (1 si le
compte détient le profil, 0 sinon). Une colonne par paire interdite~:
\begin{lstlisting}
G2 (LIQ+ORD) : =SI(ET(C2=1;D2=1);"CONFLIT";"")
H2 (FRS+ORD) : =SI(ET(E2=1;D2=1);"CONFLIT";"")
I2 (FRS+LIQ) : =SI(ET(E2=1;C2=1);"CONFLIT";"")
J2 (ADM+...) : =SI(ET(F2=1;SOMME(B2:E2)>0);"CONFLIT";"")
\end{lstlisting}
\end{boxoutils}

\begin{boxscenario}{Dossier MINÉQUIR. Les conflits}
\textbf{Constatation.} Cinq comptes réunissent LIQ et ORD, dont
\texttt{mbarga.j}~; deux réunissent FRS et ORD (\texttt{mbarga.j} et
\texttt{nkolo.r}). Aucune revue des habilitations n'a été faite depuis la
mise en service de SIDÉPENSE en 2024.

\medskip
\textbf{Analyse} (confiance élevée). Le conflit LIQ + ORD n'est pas un
cas isolé~: il touche cinq comptes, ce qui montre que la séparation n'a
jamais été conçue dans l'attribution des profils. L'absence de revue
explique que personne ne l'ait vu.
\end{boxscenario}

\separateur

% ============================================================================
\section{La revue des habilitations}
\label{sec:c09-revue}
% ============================================================================

\begin{boxdef}
\termecle{Revue des habilitations}\index{revue des habilitations}~:
contrôle détectif par lequel le propriétaire de l'application examine,
à intervalle fixé, la liste des comptes et de leurs profils, confirme ou
retire chaque droit, et signe. Elle n'est efficace que si la liste est
exhaustive et extraite de l'application, et si les retraits décidés sont
exécutés.
\end{boxdef}

\begin{boxobjet}{C-ACC-03 · Revue des habilitations}
\textbf{Objectif.} Les droits dans SIDÉPENSE restent conformes aux
fonctions.

\textbf{Risque couvert.} Accumulation de droits, comptes oubliés.

\textbf{Critère.} Revue semestrielle signée par le DAF (bonne pratique~;
aucune règle interne).

\textbf{Test de conception.} Une procédure de revue existe-t-elle~?
Non.

\textbf{Test de fonctionnement.} Sans objet~: aucune revue n'a eu lieu.

\textbf{Cote.} \NON{}.
\end{boxobjet}

\separateur

% ============================================================================
\section{L'authentification}
\label{sec:c09-authentification}
% ============================================================================

L'authentification relie un compte à une personne. Les contrôles attendus
sont simples~: un compte par personne, un mot de passe personnel, une
longueur minimale, un changement périodique ou une authentification à
deux facteurs pour les profils sensibles, un verrouillage après plusieurs
échecs. Ces règles relèvent aussi de l'audit de sécurité de l'ANTIC
(Chapitre~\ref{chap:institutions})~: le vérificateur les note sans
conduire lui-même un audit de sécurité.

\begin{boxalert}
\textbf{Violation typique.} Constater que les mots de passe n'expirent
jamais, et conclure que la sécurité de SIDÉPENSE est défaillante.

\smallskip
\textbf{Règle~:} rester dans son objectif. Pour le contrôle interne de la
dépense, le constat est que l'imputation des opérations aux personnes
n'est pas garantie~; la sécurité d'ensemble relève de l'audit de
sécurité.
\end{boxalert}

\separateur

% ============================================================================
\section{Lecture du domaine}
\label{sec:c09-lecture}
% ============================================================================

\begin{tableaucours}{L{6.2cm}cL{4.4cm}}{%
    \tableauHeader{Contrôle} &
    \tableauHeader{Cote} &
    \tableauHeader{Pièce}
}
Création sur demande validée & \NON{} & Aucune demande archivée \\
Profils sans combinaison interdite & \NON{} & Matrice (5 et 2 conflits) \\
Désactivation à la date du départ & \NON{} & 9 comptes sur 31 \\
Revue des habilitations & \NON{} & Aucune revue \\
Comptes nominatifs & \NON{} & 3 comptes génériques \\
Journalisation des actions des comptes ADM & \OUI{} & Journal \\
\end{tableaucours}
\vspace{-0.8em}
\begin{center}{\small\itshape Tableau~9.3. Lecture du domaine D1
(accès)}\end{center}

\begin{boxlecture}{Domaine D1, accès à SIDÉPENSE}
\textbf{Constatation.} Un contrôle sur six fonctionne.

\textbf{Analyse} (confiance élevée). Le domaine est au niveau~0 ou~1 de
l'échelle du Chapitre~\ref{chap:eval-sci}. La déficience est majeure~:
elle prive de valeur tous les contrôles d'application de SIDÉPENSE.

\medskip
\hfill\conformite{1}{6}{0}
\end{boxlecture}

\begin{boxmetier}{Vérificateur}
Demandez toujours trois listes extraites devant vous~: les comptes et
leurs profils, la dernière connexion de chaque compte, et les
détenteurs des profils d'administration. Croisez-les avec le fichier du
personnel du même jour.
\end{boxmetier}

\begin{boxmetier}{Contrôleur financier}
Votre visa électronique est un profil. Si un compte générique ou un autre
agent peut viser à votre place, votre contrôle n'existe pas~: signalez-le
par écrit.
\end{boxmetier}

\begin{boxmetier}{Auditeur interne}
Proposez à la direction une revue semestrielle des habilitations, signée
par le propriétaire de l'application~: c'est le contrôle le moins coûteux
du domaine, et celui qui aurait détecté les cinq conflits.
\end{boxmetier}

\separateur

\begin{boxresume}
\begin{itemize}
    \item \textbf{Cinq étapes, cinq contrôles.} Création, attribution,
          modification, revue, désactivation.
    \item \textbf{Le propriétaire décide, la cellule exécute.}
    \item \textbf{Quatre comptes à risque.} Partis, génériques, à
          privilèges, inactifs.
    \item \textbf{Croiser avec le personnel}, le même jour, sur des listes
          extraites devant l'auditeur.
    \item \textbf{La matrice révèle les conflits}~; la revue les fait
          disparaître.
    \item \textbf{Rester dans son objectif.} L'authentification se note~;
          l'audit de sécurité relève de l'ANTIC.
\end{itemize}
\end{boxresume}

\begin{boxexo}[title={\ding{45}\quad Exercice~9.1. Classer les comptes \niveaubase}]
Classez selon le Tableau~9.2~: (a) \texttt{admin}~; (b) \texttt{mbarga.j}
au 23~mars~2026~; (c) \texttt{saisie01}~; (d) un compte dont la dernière
connexion date de juin~2025.
\end{boxexo}

\begin{boxexo}[title={\ding{45}\quad Exercice~9.2. Une formule \niveaubase}]
Écrivez la formule qui signale, dans la matrice, un compte qui détient à
la fois ENG et ORD, puis dites si cette combinaison doit être interdite.
\end{boxexo}

\begin{boxexo}[title={\ding{45}\quad Exercice~9.3. La revue qui ne sert à rien \niveaumoyen}]
À la CAPNO, une revue des habilitations est signée chaque semestre. Elle
porte sur une liste saisie à la main par l'informaticien, et aucun retrait
n'a jamais été décidé. Cotez le contrôle et rédigez le constat.
\end{boxexo}

\begin{boxexo}[title={\ding{45}\quad Exercice~9.4. Un petit service \niveauexpert}]
La recette des finances de Bafoussam n'a que trois agents pour
e-RECETTES. Proposez une attribution des profils, les contrôles
compensatoires nécessaires et la preuve que chacun doit laisser.
\end{boxexo}

\begin{boxinfo}
\textbf{Transition.} Les accès maîtrisés, il faut encore que les
programmes et leurs paramètres ne changent pas sans contrôle, et que
l'exploitation tienne. C'est l'objet du Chapitre~\ref{chap:exploitation}.
\end{boxinfo}
